\documentclass[12pt,a4paper]{article}

\usepackage[utf8]{inputenc}
\usepackage[T1]{fontenc}
\usepackage{amsmath,amssymb,amsthm}
\usepackage{graphicx}
\usepackage{booktabs}
\usepackage{hyperref}
\usepackage[margin=1in]{geometry}
\usepackage{enumitem}
\usepackage{xcolor}
\usepackage{float}
\usepackage{caption}
\usepackage{subcaption}
\usepackage{longtable}
\usepackage{array}
\usepackage[nopatch=footnote]{microtype}
\emergencystretch=1.5em
% Typographic tolerance: report only boxes worse than badness 2999
% (default 1000 flags mildly-stretched lines around unbreakable \texttt
% tool names and inline math; 3000 is the conventional acceptability bound).
\hbadness=2999
\vbadness=2999
% Title page spacing via geometry
\usepackage{authblk}

\definecolor{proven}{RGB}{0,128,0}
\definecolor{assumed}{RGB}{200,150,0}
\definecolor{broken}{RGB}{200,0,0}
\definecolor{conditional}{RGB}{200,100,0}
\definecolor{dissolved}{RGB}{128,128,128}

\newcommand{\countnote}[1]{{\small\textit{(count: #1)}}}

\title{Orchestrated AI Research: Epistemic Architecture and Multi-Agent Methodology in Computational Theoretical Physics}

\author[1]{Ryan Berry}
\author[2]{Claude (Claude Code agent framework, Opus 4.6--4.8 generations)}
\affil[1]{Independent Researcher}
\affil[2]{Anthropic --- Claude Code CLI Agent Framework}

\date{Revised June 2026 (original draft March 2026)\\ \small Manuscript revision: Claude Fable 5}

\begin{document}

\maketitle

\begin{abstract}
We describe the epistemic architecture and multi-agent methodology developed over more than one hundred sessions of computational theoretical physics research using Claude Code, Anthropic's CLI agent framework. The project---an investigation of phonon-exflation cosmology on the product manifold $M^4 \times SU(3)$---served as both a physics research program and a testbed for orchestrated AI-assisted research methodology. The research has run continuously from late 2024 through the present (June 2026); this manuscript was first drafted in March 2026 at session~51 and is here substantially revised and expanded to cover the project as it now stands (session~106 and ongoing). Over this span the project deployed 34 specialist and utility AI agents across 3{,}128 computation scripts, indexed 2{,}301 theorem-class entities and 195 closed-mechanism records into a structured knowledge graph of 103{,}767 entities, and pre-registered 3{,}241 computational gates. Each of these figures is reported below with the method by which it was counted; that discipline---never citing a number without its provenance---is itself one of the methodological commitments the paper describes. We report on: (1) the epistemic discipline framework that governs evidence classification, including pre-registered gates, source authority hierarchies, and the separation of bookkeeping from reasoning; (2) the evolution of multi-agent collaboration formats from sequential workshops to parallel independent computation, now codified in a family of session skills; (3) the session-level infrastructure including carry-forward protocols, working-paper templates, and the knowledge-index/MCP pipeline; (4) the role of cross-domain research correlation in producing discoveries inaccessible to individual specialist agents; (5) a \emph{methodology-floor} layer---a maturation, over sessions~78--106, of pre-registration into a formal under-specification taxonomy, of process rules into a graded promotion contract, and of multi-agent agreement into a context-free cross-axis verification protocol; and (6) a quantitative, deliberately honest account of methodology effectiveness. On the framework's own probability: the last \emph{calibrated} estimate was 2--4\% at session~51; after session~66 the project deliberately ceased quoting a single probability, judging that post-S66 evidence is better described as the geometry of a constraint surface than as a scalar, and a formal re-anchoring remains a pending task. We treat this refusal-to-quote as itself a methodological result. We argue that the methodology's principal innovation is the \emph{epistemic architecture}: a layered system of rules, formats, audit invariants, and institutional-memory structures that allows an ensemble of memoryless AI agents to conduct rigorous falsification-driven research while maintaining coherence across sessions.
\end{abstract}

\tableofcontents

% =====================================================================
\section{Introduction}
\label{sec:intro}
% =====================================================================

The application of large language models (LLMs) to scientific research has progressed rapidly from literature summarization and code generation to substantive contributions in mathematics~\cite{alphageometry}, protein structure prediction~\cite{alphafold}, and materials discovery~\cite{gnome}. However, most deployments treat AI as a tool---a single agent executing a well-defined task. The question of whether AI systems can participate in \emph{open-ended theoretical research}---where the question itself evolves, negative results reshape the investigation, and the research program spans months or years---remains largely unexplored.

This paper reports on a computational theoretical physics project, ongoing since late 2024, that developed, tested, and iteratively refined a methodology for orchestrated multi-agent AI research. The physics context---an investigation of whether the Standard Model particle spectrum and cosmological observables can emerge from the spectral geometry of a Kaluza--Klein compactification on $SU(3)$---is documented in a companion manuscript~\cite{framework_paper}. Here we focus on the \emph{methodology}: the epistemic rules, collaboration formats, audit infrastructure, and institutional memory that have enabled coherent investigation across more than one hundred sessions.

\paragraph{A note on this revision.} This manuscript was first drafted in March 2026, when the project stood at session~51 and had run for roughly ten weeks. That draft is here substantially revised. The project did not end at session~51; it continued through session~106 (the current frozen-and-stamped state at the time of revision) and is ongoing. The original draft was, in retrospect, an accurate account of the project's \emph{first third}. Two things motivate the revision. First, every quantitative claim in the original was a session-51 snapshot, and most are now off by integer multiples (one, the count of extracted equations, must be corrected \emph{downward}---see Section~\ref{sec:assessment}). Second, and more importantly, the project's most defensible methodological contribution---a \emph{methodology floor} of formal pre-registration discipline, graded process rules, and structurally-independent multi-agent verification---did not exist when the original was written. Section~\ref{sec:floor} is devoted to it.

\subsection{The Methodological Challenge}

AI-assisted research faces a fundamental tension: LLM agents have no persistent memory across sessions, yet theoretical research requires sustained investigation in which each session builds on all prior results. A single session may produce a theorem that constrains all future work, close a mechanism that was the leading candidate, or retract a prior claim that downstream computations relied upon. Without institutional memory, each session starts from zero.

The standard solutions, context windows and retrieval-augmented generation, address the technical problem of information \emph{access} but not the epistemic problem of \emph{what to trust}. A knowledge base that treats all prior results equally cannot distinguish a machine-epsilon-verified theorem from a preliminary estimate that was later retracted. The methodology described here addresses this through what we call the \emph{epistemic architecture}: a layered system of rules, classifications, and authority hierarchies that allow agents to navigate an accumulated, partly self-contradicting research record with appropriate confidence calibration.

\subsection{Contributions}

\begin{enumerate}[label=(\arabic*)]
    \item We describe the \textbf{epistemic discipline framework} (Section~\ref{sec:epistemic}) governing evidence classification, including pre-registered gates, the constraint methodology, and the source authority hierarchy that ranks adversarial skeptic verdicts above synthesis files above raw computation output.

    \item We document the \textbf{evolution of collaboration formats} (Section~\ref{sec:formats}) from early sequential workshops through panel debates to parallel independent computation, now codified in a family of session-execution skills.

    \item We present the \textbf{session infrastructure} (Section~\ref{sec:infrastructure}): the carry-forward protocol, the working-paper template that resolved multi-agent file contention, and the knowledge-index / MCP pipeline.

    \item We analyze the \textbf{role of cross-domain correlation} (Section~\ref{sec:crossdomain}) in producing discoveries no individual specialist agent could make, with a case study and its later generalization into a typed bridge formalism.

    \item We describe the \textbf{methodology floor} (Section~\ref{sec:floor}): the apparatus that matured over sessions~78--106, including a pre-registration under-specification taxonomy, a graded SUGGESTION$\to$MANDATORY rule-promotion contract, and a context-free cross-axis verification protocol that is the constructive answer to ``agreement among agents is not evidence.''

    \item We provide a \textbf{quantitative, deliberately honest methodology assessment} (Section~\ref{sec:assessment}), including a full statistic recount with counting methods, a retraction analysis, and the framework's evolving (and partly suspended) probability-reporting methodology.
\end{enumerate}

% =====================================================================
\section{Epistemic Discipline Framework}
\label{sec:epistemic}
% =====================================================================

The epistemic framework is codified in rule documents that every agent loads automatically at spawn time. These are not guidelines---they are enforced constraints; violations trigger immediate correction by the human or by an automated audit.

\subsection{The Constraint Methodology}

The foundational principle is \emph{pre-registered falsification}:

\begin{quote}
Pre-register gates BEFORE computation---define pass/fail criteria, then compute. Report the number first. Classify second. Interpret third.
\end{quote}

Every computation in the project has a pre-registered gate with explicit PASS, FAIL, or INFO criteria defined before the computation runs. This eliminates the most common failure mode of AI-assisted research: post-hoc rationalization of an ambiguous result. The principle is simple to state and was easy to honor at session~51; the bulk of the project's later methodological work (Section~\ref{sec:floor}) consists of recognizing that pre-registering a gate's \emph{pass criterion} is insufficient unless the \emph{machinery} the gate depends on is also pinned in advance.

A critical corollary: \emph{negative results are boundaries, not failures}. Each closure constrains the solution space, narrowing the region where viable mechanisms can exist. A framework that has tested and closed many wrong mechanisms is stronger than one that has tested none; eliminating a wrong corridor is information.

\subsection{Source Authority Hierarchy}

When sources conflict, higher authority wins:

\begin{enumerate}
    \item Adversarial-skeptic verdicts (highest)
    \item Synthesis files
    \item Gate verdict results
    \item Session minutes
    \item Raw computation output (lowest)
\end{enumerate}

This hierarchy encodes a deliberate principle: \emph{the agent whose role is to say ``no'' has the highest authority}. A dedicated empirical-skeptic agent---whose function is adversarial skepticism and which is the project's sole estimator of framework probability---outranks all specialist assessments. This counteracts the natural tendency of a specialist agent to over-value a result in its own domain.

\subsection{Confidence Calibration}

The framework prohibits filler confidence language (``promising,'' ``encouraging,'' ``likely correct'') in all agent output. Pre-registered gates are the only evidence; everything else is commentary. This seemingly minor rule had outsized impact: it forced agents to commit to specific numerical thresholds rather than hedging qualitatively. The probability trajectory (Section~\ref{sec:assessment}) reflects this discipline---the framework probability dropped sharply in single sessions when decisive gates fired, rather than drifting through accumulated soft language.

\subsection{The Assumptions Registry and the DISSOLVED Category}
\label{sec:dissolved}

Session~50 produced a formal \emph{Foundation and Assumptions Registry} classifying every assumption in the framework. It remains live, maintained as a curated atlas document. As of the most recent stamp it contains 63 entries \countnote{assumptions-registry atlas header; entry table} in five status classes (a sixth, \texttt{STAGE-1-CANDIDATE}, is under consideration):

\begin{itemize}
    \item \textcolor{proven}{\textbf{PROVEN}} (20): Derived from mathematics, verified to machine epsilon
    \item \textcolor{assumed}{\textbf{ASSUMED}} (15): Input to the framework, not derived
    \item \textcolor{conditional}{\textbf{CONDITIONAL}} (12): Depends on uncomputed or scheme-dependent quantities
    \item \textcolor{broken}{\textbf{BROKEN}} (11): Tested and failed
    \item \textcolor{dissolved}{\textbf{DISSOLVED}} (5): Question reframed; assumption no longer relevant
\end{itemize}

\noindent\textit{(Counting method: tallies read from the assumptions-registry atlas header, current stamp; the original March draft reported the session-50 snapshot of 15/14/8/11/4.)}

The DISSOLVED category proved essential for preventing false dichotomies. When the spectral-action stabilization mechanism was proven impossible (a structural monotonicity theorem, session~37), the question ``What potential well stabilizes the modulus at the fold?'' did not become BROKEN---it was DISSOLVED, replaced by a transit paradigm in which the modulus rolls through the fold without stopping. The assumption was not wrong; it was the wrong question. The two dozen equilibrium-stabilization closures from sessions~17--40 are preserved in the registry as DISSOLVED via this transit reframe rather than counted as failures---a concrete instance of the constraint methodology distinguishing ``a corridor is closed'' from ``the program failed.''

% =====================================================================
\section{Evolution of Collaboration Formats}
\label{sec:formats}
% =====================================================================

The project tested several distinct multi-agent formats, each encoding different assumptions about how specialist knowledge combines. The qualitative lessons below were learned over sessions~16--51 and remain canonical; the formats themselves have since been consolidated into a family of reusable execution skills (Section~\ref{sec:rclab}).

\subsection{Workshop Format}

Sequential paired discussions in rounds. Two agents discuss a topic, produce a markdown document, and the next pair reads all prior output. This format excels at deep cross-domain exploration but scales poorly: each round waits for the previous round's output, and total session time grows linearly with agent count.

\textbf{Failure mode}: agent \emph{roleplaying}---when one agent writes another's workshop section, lacking that agent's domain knowledge, producing confident but shallow text. This was identified as a recurring anti-pattern and prohibited by rule: an agent must never write another agent's designated section.

\subsection{Panel Format}

Interpretive panel with a designated writer. Two or three specialists interpret shared material from their domain perspective; a designated writer synthesizes. This produces the richest single-document output but requires careful writer selection and is vulnerable to the ``agents lie about being done'' phenomenon: an agent claims completion but produces its best cross-talk \emph{after} the initial claim. The operative rule that emerged: \emph{only the human confirms completion}; an agent's self-report of ``final'' is not evidence that it is finished.

\subsection{Parallel Independent Computation}

Each computation is an independent agent call with a self-contained prompt. No team, no inbox, no inter-agent messaging. Agents write to designated sections of a shared working paper; decision points between waves are evaluated by a coordinating agent (the ``team lead'') or the human.

\textbf{Why it works}: it eliminates the notification avalanche that plagues team-based formats (four or more agents in a team produce $O(n^2)$ cross-notifications), eliminates inbox contamination between concurrent teams, and allows unlimited parallelism within each wave. The working paper provides shared state without requiring real-time communication. The critical innovation was the \emph{working-paper template}: pre-labeled sections, one per agent per round, which solve file contention, enable cross-reading, and provide an audit trail. As the project's own notes put it, discovering a markdown template that does this beats elaborate inter-agent message choreography.

\subsection{The rclab Skill Family (current pipeline)}
\label{sec:rclab}

The three formats above evolved into a family of eight reusable session-execution skills \countnote{\texttt{ls .claude/skills/rclab-*}}, which now constitute the standard pipeline:

\begin{itemize}
    \item \texttt{rclab-plan} --- plans the next compute session: gathers carry-forwards from prior working papers, partitions them into waves by theme, spawns a per-wave planner swarm that writes full-fidelity gate specifications, validates upstream dependencies, and maintains the forward-priority registers.
    \item \texttt{rclab-coordinate} --- executes a session plan in compute mode: hands each test case to its agent, waits, reports. This is the direct successor of the parallel-independent-computation format; no teams, no inboxes.
    \item \texttt{rclab-investigate} --- generates a workshop-schedule campaign from a just-closed session, identifying genuine adversarial workshops (Section~\ref{sec:floor-workshop}) and routing non-workshop items to carry-forward.
    \item \texttt{rclab-review} --- independent solo synthesis: several agents each read the same sources and write their own report, with no coordination between them.
    \item \texttt{rclab-workshop} --- a two-agent iterative workshop on a shared document, sequential, with no team infrastructure.
    \item \texttt{rclab-team} --- a coordinated multi-agent team with named-agent inbox messaging, for panel and multi-round interpretive formats.
    \item \texttt{rclab-solo} --- single-agent sequential wave execution.
    \item \texttt{rclab-reflect} --- structured reflection on a just-closed wave, plan, or document: what stood out, what was surprising, what crosses gates.
\end{itemize}

The original draft described an earlier two-stage ``Collab-Plan / Collab-Team'' pipeline; the rclab family is its descendant, split into finer-grained skills as the distinction between independent computation, adversarial workshops, and independent review hardened into separate routing decisions (Section~\ref{sec:floor-workshop}).

% =====================================================================
\section{Session Infrastructure}
\label{sec:infrastructure}
% =====================================================================

\subsection{The Carry-Forward Protocol}

Every session produces recommendations: computations suggested by reviewers, open questions, pre-registerable predictions. The carry-forward rule requires that session plans carry forward all genuine future-computation recommendations as planned computations; an item not planned is, in practice, lost. This rule emerged from an early session in which several prior recommendations were silently dropped, and the computations they would have driven were never performed.

The protocol has since been sharpened in two directions that the original draft predates. First, a \emph{fix-in-session} discipline: a deviation that can be corrected now (a status-tag edit, a hygiene repair) is fixed now, not logged and carried forward. Second, a \emph{no-padding} rule: a carry-forward is reserved for genuine future computation and must carry a four-field specification (what / inputs / gate / effort); hygiene observations on already-correct artifacts are explicitly \emph{not} carry-forwards. The two rules together produce a shorter, signal-rich forward queue rather than an inflated one, and the distinction between in-session resolution and genuine carry-forward is enforced through per-session housekeeping ledgers (Section~\ref{sec:floor-workshop}).

\subsection{The Knowledge Index and Query Infrastructure}
\label{sec:knowledge}

The project maintains a structured knowledge graph that has grown into a multi-component query infrastructure. The original draft described this as a single 30~MB JSON file; it is now the following stack \countnote{methods inline per component}:

\begin{itemize}
    \item \texttt{knowledge-index.json} --- the canonical knowledge graph, approximately 45~MB and 1.22~million lines \countnote{\texttt{ls -la}; \texttt{wc -l}}, with nine entity types (theorems, closed mechanisms, gates, sessions, open channels, researchers, equations, data-provenance records, and constants).
    \item \texttt{knowledge.db} --- a SQLite database with a full-text (FTS5) search index, rebuilt from the JSON graph. It holds 103{,}767 full-text-indexed entities \countnote{\texttt{COUNT(*)} on the FTS table}, including 19{,}713 equations, 71{,}102 inter-entity edges, and 2{,}941 data-provenance records.
    \item A dedicated \textbf{knowledge MCP server} exposing the database to agents through a small tool set: full-text-ranked \texttt{search\_knowledge}, \texttt{query\_entity}, \texttt{trace\_entity} (evidence chains), \texttt{get\_constant} (value plus provenance), and a race-safe \texttt{emit\_verdict} writer (Section~\ref{sec:floor-verdict}). At the time of revision the server's usage counter recorded 11{,}637 calls \countnote{server usage counter}, dominated by \texttt{search\_knowledge} (5{,}634) and \texttt{get\_constant} (3{,}624)---evidence of how heavily agents lean on querying the canonical graph rather than recomputing.
    \item A \texttt{routing\_manifest.json} export consumed by a separate static-site build, and a \texttt{/weave} maintenance pipeline (extractor $\to$ audits $\to$ database sync $\to$ manifest regeneration).
\end{itemize}

A standing rule requires every computation agent to query this graph \emph{before} computing, to check whether a result is already known, a mechanism already closed, or a gate already evaluated. This ``query first, compute second'' discipline exists because agents repeatedly rediscovered settled results; the knowledge infrastructure exists to make rediscovery detectable and avoidable.

\subsection{Agent Memory and Institutional Knowledge}
\label{sec:memory}

Each specialist agent maintains persistent memory files that survive across sessions, encoding domain-specific lessons it re-loads at spawn. The project-level memory encodes framework-wide state---proven theorems, closed mechanisms, orchestration lessons, and the human's critical feedback---and is loaded into every conversation.

The original draft treated agent memory as straightforwardly beneficial. Subsequent experience produced a sharp \emph{memory-scope discipline} (Section~\ref{sec:floor-amri}): agent memory is for agent-\emph{private} context only and is never the canonical home for data that other agents or gates cite. Project-level registries (watchlists, detector rosters, canonical-constant values, gate-verdict tables) live in queryable registry files and the knowledge database, not in memory. The reasoning is traceability: a value other gates depend on must live where the audit trail can reach it.

% =====================================================================
\section{Cross-Domain Research Correlation}
\label{sec:crossdomain}
% =====================================================================

The project's most significant early methodological finding is that \emph{discoveries at domain intersections require a generalist coordinator with access to the full research corpus}, not just specialist agents operating within their domains.

\subsection{Case Study: The Spectral Action Correlator}
\label{sec:sacase}

In one session the project deployed eight specialist agents to test escape routes from an exact algebraic identity, $\alpha_s = n_s^2 - 1$, that threatened a cosmological prediction. All eight tested variations of the same Josephson-phase correlator family. All eight failed---the identity is structural within that family.

The breakthrough came when the coordinating agent (not a specialist) read three papers from different research domains: a spectral-geometry result on the propagator structure of the spectral action; a prior internal review noting that the spectral action and the Josephson coupling are structurally different functionals that should not be mixed; and a superfluid-physics result that, in non-equilibrium superfluids, different correlators (NMR, sound, superfluid density) give different momentum dependence. The connection: the spectral-action two-point function is a different mathematical object from the Josephson propagator---a much broader pole spread, not protected by Goldstone's theorem---and so it \emph{does} break the identity. No individual specialist agent could have made this connection, because it required simultaneously understanding noncommutative-geometry spectral-action structure, condensed-matter multi-band correlators, and superfluid multi-probe phenomenology.

\subsection{From Ad-Hoc Connection to Typed Bridge}

This kind of cross-domain connection was later generalized into a formal device: the \emph{cross-pillar bridge anatomy} (Section~\ref{sec:floor-bridge}). Rather than relying on a coordinator's lucky association, the framework now requires that any claimed connection between an observable on one domain (``pillar'') and a measurement on another declare five explicit anatomy elements and a three-level confidence ladder. The ad-hoc S50 connection became the prototype for a typed, auditable bridge formalism.

\subsection{Quantifying Cross-Domain Discovery (with a caution)}

The project's Breakthrough Genealogy catalogues 39 major breakthroughs \countnote{breakthrough-genealogy atlas header}; the original draft reported 15. The catalogue spans structural theorems, wall discoveries, paradigm shifts, observational matches at zero free parameters, and---in two categories that postdate the original draft---methodology-floor advances and calibration-corpus (rule-promotion) events.

\textbf{An honesty caution, reproduced from the catalogue itself.} Of the 39 catalogued breakthroughs, ten ($25.6\%$) are methodology-floor or calibration-corpus events. These advance \emph{process discipline}, not the framework's substrate-physics probability. A breakthrough that promotes a rule from SUGGESTION to MANDATORY makes the framework's pre-registration more rigorous; it does \emph{not} make the framework more likely to be true about the physics. Counting methodology-floor breakthroughs alongside substrate-physics breakthroughs in any directional or probability assessment is an inflation error the catalogue explicitly flags and that this paper, being about honest methodology, reproduces rather than launders. The genuinely cross-domain \emph{substrate-physics} breakthroughs (block-diagonality from NCG $\times$ condensed-matter sector decomposition; the BCS chain from condensed-matter $\times$ nuclear structure; the instanton-gas paradigm from nuclear backbending $\times$ integrability; the spectral-action correlator from NCG $\times$ superfluid $\times$ resonance) remain the supporting evidence for the cross-domain thesis.

% =====================================================================
\section{The Methodology Floor (Sessions 78--106)}
\label{sec:floor}
% =====================================================================

This section describes the apparatus that did not exist when the original draft was written and that now constitutes the project's most defensible methodological contribution. We call it the \emph{methodology floor}: the infrastructure that governs \emph{how} the substrate is investigated, as distinct from \emph{what} the substrate is. It is curated as a dedicated synthesis document; this section renders that synthesis for a reader who has never seen the project, defining terms as it goes.

A single organizing idea recurs: a \emph{layer functor} $F$ mapping a substrate-physics quantity to its methodology-floor image to its audit-floor image. Under $F$, an eigenvalue maps to rule-file content, a numerical PASS predicate maps to an artifact-existence predicate, and a verdict's numerical value maps to a cryptographic hash of the verdict's inputs. The point of the functor is that the same structural failure---for instance, under-specifying something that should have been pinned in advance---recurs identically at the physics layer (an unpinned machinery parameter), the methodology layer (an unpinned rule criterion), and the audit layer (an unpinned audit input). Recognizing the shared structure let the project build each floor by analogy with the one above it.

\subsection{Pre-Registration Maturation: PRU, PRDR, and Source Reconciliation}
\label{sec:floor-pru}

The original draft's ``pre-register the gate'' principle proved insufficient. A pass criterion can be fixed in advance while the \emph{machinery} the gate depends on (a regularization scheme, a truncation order, a representation convention) is left free, creating execution-time latitude indistinguishable from post-hoc tuning. The response was a formal taxonomy.

\textbf{PRU (Pre-Registration Underspecification)} classifies the ways a plan can leave a gate-relevant machinery parameter unpinned. Its principal sub-classes (collectively ``Class 8'') are: machinery-pin cardinality failure; verifier-rubric pre-registration failure (when a PASS criterion involves grading qualitative content, the grading rubric itself must be pre-registered, or execution-time rubric tuning is indistinguishable from iterate-until-PASS); output-precision pre-registration failure (a value cited downstream must declare its significant figures, so a downstream verifier's tolerance cannot be set tighter than the value's own precision); and further sub-classes for representation-convention, joint-hypersurface form, layered substitution chains, and degenerate observables.

\textbf{PRDR (Pre-Registration Dry-Run)} is the preventive procedure: before a plan is frozen, the producing script is dry-run, every free parameter is enumerated by static analysis, and each is either pinned or explicitly declared diagnostic.

\textbf{Source Reconciliation} is a separate, commuting audit that detects pinned-but-\emph{drifted} values (as opposed to PRU's missing values). It compares each pin against the canonical value in the knowledge database and classifies the drift into six classes (for example: pin-band tighter than canonical; pin-band looser than canonical, the highest-leverage false-pass direction; pin computed against a since-superseded canonical; pin given as a placeholder when a substrate-first canonical exists). Severity scales with the order-of-magnitude discrepancy, escalating to a hard halt of the plan freeze.

Together these turned ``pre-registration'' from a one-sentence principle into a multi-stage pipeline run at every plan freeze: cardinality pre-flight, then value drift, then source existence, then machinery enumeration, then execution.

\subsection{Graded Process Rules: the K-Counter Promotion Contract}
\label{sec:floor-kcounter}

Methodology rules are not born mandatory. A new rule enters as a \textbf{SUGGESTION} and is promoted to \textbf{MANDATORY} only after accumulating three \emph{structurally distinct} calibration instances (denoted $K=3$)---three genuinely different occasions on which the rule caught a real problem, not three repetitions of the same case. The calibration instances are recorded in an \emph{append-only, forward-only} corpus: instances are added, never removed or reordered, so the promotion history is itself an audit trail.

This is, to our knowledge, a novel piece of epistemic infrastructure: \emph{graded, evidence-accumulating process rules}. It addresses two failure modes symmetrically. A rule imposed by fiat with no calibration tends to be over-broad and to ossify; a rule never promoted stays advisory and is ignored. The K-counter forces a rule to earn its mandatory status on accumulated evidence, and it is governed by a meta-rule whose content is that one should replace lists of mandatory imperatives with examples, schemas, and pointer tables---structure that compensates for what a flat rule list cannot express.

\subsection{Multi-Agent Agreement Done Right: Context-Free Cross-Axis Verification}
\label{sec:floor-stage2}

The original draft implicitly relied on agent agreement as a positive signal while never quite confronting the problem that agreement among agents sharing context is not independent confirmation---shared context produces shared output. The methodology floor supplies the constructive answer: a four-stage promotion pathway for any \emph{joint} theorem (one whose statement requires evidence from more than one methodological axis).

\begin{description}
    \item[Stage 0 --- Workshop-internal candidate.] The theorem is drafted inside a workshop by the authoring agents, with each clause attributed to the axis that authored it. The workshop verdict freezes the text.
    \item[Stage 1 --- Registered candidate.] The frozen text is registered in the permanent results registry tagged \texttt{STAGE-1-CANDIDATE}, with its joint clauses flagged. Downstream work may cite it only with the candidate qualifier.
    \item[Stage 2 --- Context-free cross-axis verification.] Two independent reviewers, one per axis, are dispatched in parallel. Critically, they \emph{operate without prior workshop context}: they read only the registered Stage-1 entry, never the workshop transcripts, and they cannot be the original authoring agents. Each verifies the clauses on its own axis plus all joint clauses; joint clauses must pass in \emph{both} verdicts (logical AND). A reviewer-selection protocol enforces axis-distinctness, exclusion of the original authors \emph{and} of any successor agent whose memory inherits the workshop's reasoning, and audit-coverage adequacy.
    \item[Stage 3 --- Permanent.] On a both-axes pass, the tag becomes \texttt{STAGE-3-PERMANENT}.
\end{description}

The ``without prior workshop context'' condition is the whole point. If two reviewers, given only the registered statement and never the derivation, independently confirm a joint clause, their agreement is \emph{structurally} independent rather than shared-context agreement. This is the precise distinction the rule draws: agreement among agents who read the same workshop is still not evidence; agreement among agents who provably did not is. A worked instance ran at session~100a as a three-agent cross-axis verification (spectral, transit, and semiclassical reviewers; all three non-authors; clean prompts) returning nine of nine clause-verdicts in agreement, promoting the theorem to permanent.

\subsection{Enforceable Verdict Permanence: Dual-SHA, Race-Safe Emission, and Bounded Recovery}
\label{sec:floor-verdict}

The original draft asserted that ``gate verdicts are permanent.'' The methodology floor makes that enforceable. Each computation verdict is an append-only line carrying a \textbf{dual SHA-256} pair: a \texttt{content} hash over the result, and an \texttt{audit} hash over the gate's input-pin map. SHA uniqueness across all verdict lines is a checked invariant---a duplicate audit hash signals a copy-pasted or hardcoded value rather than a freshly computed closure, and is caught automatically.

Verdicts are written through a \textbf{race-safe} emission tool (a single lock-serialized writer), because raw file appends are not atomic across processes on the project's platform and concurrent agents otherwise lose lines. A \textbf{closure ladder} scores five signals over a session's verdicts (machinery-pin completeness, dual-SHA presence, completion-queue coverage, schema conformance, and SHA uniqueness). Failures route through a \textbf{bounded recovery procedure}: at most two automated remediation attempts per signal, with a termination proof, and a four-item prohibited-actions set---convention-shopping, iterate-until-PASS, post-hoc editing of a frozen pre-registration, and manually forcing a PASS---any of which immediately halts automatic recovery and escalates to the human. The bounded-iteration cap closes the iterate-until-PASS pathway \emph{by construction} rather than by exhortation.

\subsection{Honest Closure versus the Task-Complete Lie}
\label{sec:floor-closure}

A specific and recurring agent failure mode was documented: an agent appends a verdict line and then terminates, its final message \emph{claiming completion}, while a promised working-paper section is never written. We call this the ``task-complete lie''; it is not malice but a premature stop. The structural response has two parts. First, a completion-verification discipline: the orchestrator verifies promised artifacts exist on disk---by content, not by line count---before accepting an agent's completion claim. Second, a \emph{mechanical-closure discipline} for the legitimate case where a gate genuinely cannot be evaluated because an upstream prerequisite did not pass. Such a gate may be closed by an orchestrator-authored script only under explicit conditions: the upstream block must be the documented cause; the emitted verdict must be FAIL or ``pre-registration incomplete,'' never PASS; the audit hash must be per-gate distinct; the verdict value must name the blocking prerequisite; and the working-paper section must be updated in the same run. This distinguishes honest ``could not evaluate, here is why'' closure from a silently skipped task.

\subsection{Memory as Agent-Private Context: AMRI}
\label{sec:floor-amri}

The memory-scope discipline of Section~\ref{sec:memory} is enforced by an audit for \emph{Agent-Memory Registry Inversion} (AMRI): the failure in which a section of an agent's memory has quietly become a de-facto registry that other gates depend on. Three tests detect it---whether another gate pins the memory file as an audit input; whether a gate writes to a memory file as its primary registry output; and whether two agents' memories hold overlapping entries for the same observable, mechanism, or detector. Any one firing routes the section for migration to a project-level registry plus the knowledge database. The discipline applies identically to per-agent memory and to the orchestrator's own project memory; the human made the latter scope explicit on the grounds that not following it breaks traceability. Several AMRI migrations are documented. This directly upgrades the original draft's untroubled treatment of agent memory: memory is for spawn-time re-use, not for canonical cross-gate data.

\subsection{Cross-Pillar Bridge Anatomy}
\label{sec:floor-bridge}

The typed-bridge formalism promised in Section~\ref{sec:crossdomain} requires every registered connection between a substrate observable on one pillar and a laboratory observable on another to declare five \emph{anatomy} elements---the substrate-side observable; the laboratory-side observable; the explicit bridge map (a named mathematical map, not the words ``analogous to''); an algebraic convergence envelope; and an empirical anchor---together with a three-level \emph{confidence ladder}: Level~1, a regulator-invariant identity at the cohomology-class level; Level~2, a truncation-dependent convergence envelope; and Level~3, an empirical anchor at a canonical truncation, which must lie inside the Level-2 envelope for the entry to count as passing. The formalism exists to prevent ``container thinking''---treating the laboratory pillar as fundamental and the substrate as derived---by fixing the direction of explanation at the registry-entry level so that downstream citations cannot drift.

\subsection{The Workshop Discriminator and Housekeeping Ledgers}
\label{sec:floor-workshop}

The original draft used ``workshop'' loosely. The term is now precisely defined---two or more agents with genuinely competing perspectives on a specific tension, multi-round structure, and an output that is a structural verdict rather than a queued computation---and is policed by a three-question discriminator applied before anything is scheduled as a workshop: is the tension a genuine physics adjudication (then it is a workshop); is it registry-state classification or hygiene (then it is a carry-forward); is it a parallel-compute structure dressed as a panel (then it is a carry-forward). The motivating failure was investigators padding ``workshop'' lists with carry-forward equations; the discriminator exists to keep an honest, small count (a typical session yields zero to four genuine workshops, and ``no workshops'' is a valid output). Non-workshop items route to per-session \emph{housekeeping ledgers} with a fixed five-section partition (in-session resolutions, hygiene carry-forwards, parallel-compute waves, rule-extensions, and pre-compute-shell escalations), which keeps the boundary between in-session bookkeeping and genuine future computation crisp.

\subsection{Living Prioritization and a Cautionary Tale}
\label{sec:floor-evoi}

Computation priority is set by an Expected-Value-of-Information (EVOI) heuristic---each candidate computation is scored by how far the framework's standing would move on each verdict, weighted by the probability of that verdict---and recorded in a living priority table. The table carries a load-bearing honesty caveat: its EVOI values are \emph{ordinal leverage proxies, not calibrated probabilities}, and are explicitly not to be quoted as probabilities. A cautionary episode is recorded in the table's own preamble: its content sat frozen for roughly thirteen sessions while broad multi-file commits touched its bytes, so it \emph{looked} maintained while no priority in it was actually refreshed. The root cause was that no planning skill consumed the table, so an unenforced rule rotted silently; the fix was enforcement wiring---a staleness audit keyed to a machine-readable content-currency tag, plus integration into the planning skill---so that the content date can no longer lag the session unseen. ``Git-touched is not content-refreshed'' became a standing lesson.

\subsection{Substrate-First Framing as Auditable Pre-Registration}
\label{sec:floor-framing}

Two framing disciplines are treated not as stylistic preference but as epistemic infrastructure. The first governs the \emph{direction} of explanation: every explanation must flow from the substrate toward emergent physics, never the reverse, with an explicit table of wrong (container-thinking) versus right (substrate-first) phrasings. The second governs \emph{sourcing}: a numerical pin must originate from a substrate-first computation, never from an external paper's value treated as authoritative, with a clear line between a \emph{methodological} citation (a cross-check or heritage reference, allowed) and a \emph{canonical} citation (supplying a number the computation should have produced, forbidden). The reason these count as methodology rather than style is that they are enforced at the verifier-rubric and plan-freeze layers: the framing rule supplies the negative-marker set against which qualitative content is graded, which makes ``substrate-first'' an auditable predicate rather than an aspiration.

% =====================================================================
\section{Quantitative Methodology Assessment}
\label{sec:assessment}
% =====================================================================

This section reports the project's quantitative state. Because the paper is about honest methodology, every figure is given with the method by which it was counted, and the original March-2026 (session-51) value is shown alongside for contrast. We deliberately avoid session-aggregate vanity metrics (for example, PASS/FAIL ratios presented as a quality score); where the original draft used such metrics, they are scoped below strictly as historical session-1--51 observations.

\subsection{Statistic Recount}
\label{sec:recount}

The authoritative source for entity counts is the knowledge database (the SQLite store the knowledge MCP serves); file counts are from the working tree.

\begin{longtable}{>{\raggedright\arraybackslash}p{0.25\textwidth} >{\raggedright\arraybackslash}p{0.13\textwidth} >{\raggedright\arraybackslash}p{0.15\textwidth} >{\raggedright\arraybackslash}p{0.36\textwidth}}
\caption{Statistic recount: original (S51) versus current, with counting method.}\label{tab:recount}\\
\toprule
\textbf{Claim} & \textbf{S51 (orig.)} & \textbf{Current} & \textbf{Counting method / source} \\
\midrule
\endfirsthead
\multicolumn{4}{c}{\tablename\ \thetable\ -- continued from previous page}\\
\toprule
\textbf{Claim} & \textbf{S51 (orig.)} & \textbf{Current} & \textbf{Counting method / source} \\
\midrule
\endhead
\bottomrule
\endfoot
Sessions & 51 & $\geq 106$ & EVOI table content-currency tag (S106); session directories exist through \texttt{session-106}. See note~(a). \\
Duration & $\sim$10 weeks & $\sim$17 months & Late 2024 -- June 2026; the ``10 weeks'' was real for S1--S51. \\
Agents (specialist + utility) & 29 & 34 & \texttt{ls .claude/agents/*.md} $=34$; of these $\approx 31$ are physics specialists, 3 are utility (coordinator, knowledge-weaver, web-researcher). \\
Computation scripts & 630 & 3{,}128 & \texttt{find computations -name "*.py" | wc -l}. \\
Theorem-class entities & 392 & 2{,}301 & \texttt{COUNT(*)} on the \texttt{theorems} table. See note~(c). \\
Closed-mechanism records & 58 & 195 & \texttt{COUNT(*)} on the \texttt{closed\_mechanisms} table; deduplicated from 735 at S90. See note~(d). \\
Gates registered & 240 & 3{,}241 & \texttt{COUNT(*)} on the \texttt{gates} table. \\
Knowledge entities & 82{,}000+ & 103{,}767 & \texttt{COUNT(*)} on the full-text index table. \\
Knowledge-index size & 30~MB & $\sim$45~MB / 1.22M lines & \texttt{ls -la}; \texttt{wc -l} on \texttt{knowledge-index.json}. \\
Equations extracted & 80{,}497 & 19{,}713 & \texttt{COUNT(*)} on the \texttt{equations} table. \emph{Lower}; see note~(c). \\
Data-provenance records & 567 & 2{,}941 & \texttt{COUNT(*)} on the \texttt{data\_provenance} table. \\
Probability-trajectory points & 172 & 212 & \texttt{COUNT(*)} on the \texttt{probability\_trajectory} table. \\
Retractions / corrections & 25 & 49 & Retraction-log atlas master total (through S102). See note~(e). \\
Major breakthroughs & 15 & 39 & Breakthrough-genealogy atlas header. See note~(f). \\
\end{longtable}

\paragraph{Note (a) --- session count.} There is no single integer ``session count.'' Three defensible figures exist: at least 106 distinct session directories with closed working papers; 123 rows in the database's \texttt{sessions} table (which over-counts because it includes sub-sessions such as 21a/21b/22b and deliberately split sessions such as 73a/73b and 100a/100b); and a count of 36 ``final'' summary files, which is \emph{not} a session count at all (that window spans only S51--S85---earlier finals were folded into the project atlas and later ones stopped being written to that directory). The honest figure is ``more than one hundred sessions,'' with the sub-session ambiguity stated. The S100a/S100b split is illustrative: a single planned session~100 was divided into a compute campaign (100a) and an adjudication campaign (100b), each with its own working papers and verdicts.

\paragraph{Note (c) --- theorem and equation counts.} The 2{,}301 figure is the extractor's harvest of theorem-\emph{class} entities from session text; it is \emph{not} 2{,}301 machine-epsilon structural proofs. The paper keeps two numbers cleanly separated: the curated set of publishable mathematical results (a handful of standalone manuscripts---see Section~\ref{sec:math-survive}) and the indexed theorem-entity count. The equation count, notably, went \emph{down} ($80{,}497 \to 19{,}713$): the original figure counted raw LaTeX-fragment hits, while the current table holds de-duplicated equation entities. Reporting only the larger number would be inflation; reporting only the smaller would hide the index scale. We report both, labelled.

\paragraph{Note (d) --- closed-mechanism count is a contested metric.} The framework's own rules state that closures cluster by topic: four agents hitting the same truncation wall are one methodological finding, and three failing mechanisms of one kind are one open problem with three eliminated approaches. The deduplication at session~90 (735 records reduced to 266) was the framework acting on exactly this principle; continued consolidation has brought the live extractor count to 195. We therefore report 195 closed-mechanism \emph{records}, deduplicated from 735, with the explicit caveat that these cluster into a substantially smaller number of independent structural findings. We do \emph{not} present a raw closure count as a quality metric; that would violate the no-vanity-metrics discipline.

\paragraph{Note (e) --- retractions.} The retraction log records 49 retractions and corrections through session~102, in a richer taxonomy than the original draft's four categories: honest computation error, overclaiming, wrong observational mapping, conceptual reframe, and---postdating the original---\emph{supersession} (document-level authority migrating to a successor while the underlying claims survive) and \emph{interpretive-degree-of-freedom rescoping} (a falsifier \emph{relocating} to a different instrument or scale rather than vanishing; for instance, a gravitational-wave discriminator whose peak frequency lay far above any detector band was re-anchored to a large-scale-structure observable). The richness is the point: a long, well-categorized retraction log is evidence of a working falsification discipline, not of failure.

\paragraph{Note (f) --- breakthroughs and the layer caution.} Of 39 catalogued breakthroughs, ten advance methodology-floor or rule-promotion process only and do not raise substrate-physics probability (Section~\ref{sec:crossdomain}). The caution is reproduced from the catalogue and must travel with the number.

\subsection{Retraction Analysis}

The retraction log's most consequential category is wrong observational mapping---computing a correct mathematical quantity and then incorrectly identifying it with a physical observable. The starkest instance reversed an apparent strong positive into a decisive negative when a model prediction that had been compared against a \emph{derived} fit parameter was instead compared against \emph{raw} data and found to be overwhelmingly excluded. This single failure mode---comparing against derived rather than raw quantities---is the subtlest and most expensive in the log, and several rules now target it directly. The qualitative trend the original draft reported (retraction frequency falling as structural walls accumulate to prevent future agents from repeating closed mistakes) remains the right reading, though we no longer attach a per-session rate to it (Section~\ref{sec:no-vanity}).

\subsection{The Mathematical Results Survive Regardless}
\label{sec:math-survive}

A consistent finding across the project's probability history is that the mathematical results survive at full confidence regardless of the cosmological fate. Several standalone mathematical manuscripts have been drafted (on the spectral index for Ornstein--Zernike operators; on a spectral-action monotonicity theorem; on a Weyl-tensor classification; on an Anderson--Higgs impossibility result; among others). These are machine-epsilon structural results about the spectral geometry of $SU(3)$; they hold whether or not the framework's cosmological predictions do. This is the appropriate granularity for the ``publishable results'' claim---a small curated set of standalone theorems---and it should not be conflated with the 2{,}301 indexed theorem-class entities of Table~\ref{tab:recount}.

\subsection{The Probability Trajectory and Its Deliberate Suspension}
\label{sec:trajectory}

The framework's probability history is the paper's most delicate honest-reporting obligation, and the place where the original draft, through no fault at the time of writing, is now most misleading.

\textbf{The calibrated history (S1--S66).} The framework's probability, as estimated by the adversarial-skeptic agent, rose from a 2--5\% prior to a 45--52\% peak at session~19d on accumulating structural mathematical results, then collapsed to roughly 8\% / 5\% at session~23a---the ``Venus Moment,'' the most decisive single-session downward revision, when a pre-registered boundary computation failed unambiguously. It recovered to 18--32\% over sessions~33--35 when non-perturbative condensed-matter physics was demonstrated, and settled to 2--4\% by session~51 as the leading cosmological predictions failed observational tests. A later phase (sessions~61--66) produced the first viable spectral-index prediction and the first viable cosmological-constant mechanism while exposing a spectral-functional sign ambiguity. The trajectory is double-peaked and moves in \emph{both} directions, which is the minimum signature of a methodology responsive to evidence rather than to confirmation bias.

\textbf{The deliberate suspension (S67 onward).} After session~66 the project \emph{stopped quoting a single calibrated probability}. The reasoning, recorded as a framework-hygiene rule, is twofold. First, the project judged that post-S66 evidence is better described as the geometry of a constraint surface---which walls are respected, which gates passed, which remain uncomputed---than as a scalar; a probability collapses that geometry into one number that invites over-reading. Second, the project introduced a structural distinction (the layer functor of Section~\ref{sec:floor}) between substrate-physics directional movements and methodology-floor movements, and recognized that the bulk of post-S66 activity is methodology-floor, which must \emph{not} be counted toward substrate-physics probability. Post-S66 sessions therefore record only \emph{directional} language (up-tick, down-tick, flat, paradigm-shift, neutral) with an explicit substrate-physics-versus-methodology-floor tag, and no number. A formal re-anchoring---a fresh adversarial-skeptic probability estimate weighting the post-S66 evidence---was recommended for session~89 and \emph{remains an open, unexecuted task} at the time of this revision.

\textbf{The honest statement, therefore}, is this: \emph{the last calibrated framework probability was 2--4\% at session~51; the project subsequently and deliberately ceased producing a single probability; a formal re-anchoring is pending.} The original draft's abstract reported ``a current 2--4\%''; that figure is the session-51 value and is, at the time of this revision, more than fifty sessions stale---and stale by a deliberate methodological choice rather than by neglect. Quoting it as current would be precisely the kind of false precision the project's own honesty discipline forbids, which is why this revision reports the suspension as a result in its own right. A methodology that knows when \emph{not} to quote a number is exhibiting the discipline this paper is about.

\subsection{On Session-Aggregate Metrics}
\label{sec:no-vanity}

The original draft included a per-session retraction rate, a closure ``efficiency'' fraction, and an agent-count-versus-discovery-rate table. The project's current rules treat such session-aggregate ratios as vanity metrics---a PASS/FAIL ratio is not a quality score, and four agents failing on one wall is one finding, not four---and forbid presenting them as live quality measures. We therefore retain these only as bounded \emph{historical} (session-1--51) observations: they describe how the early project behaved and are not extended forward as ratio-metrics, nor offered as evidence of methodological quality. The qualitative findings they motivated---that coordination overhead dominates beyond roughly four agents in a messaging team, and that the parallel-computation format eliminates that overhead---survive as design lessons (Section~\ref{sec:formats}) independent of the discarded ratios.

% =====================================================================
\section{Lessons Learned}
\label{sec:lessons}
% =====================================================================

\subsection{What Worked}

\begin{enumerate}
    \item \textbf{Pre-registered gates}, hardened into the PRU/PRDR pipeline (Section~\ref{sec:floor-pru}). Forcing commitment before computation, and forcing the \emph{machinery} to be pinned and not merely the criterion, is the single most important epistemic innovation.
    \item \textbf{Parallel independent computation} with a working-paper coordination layer (Section~\ref{sec:rclab}). Eliminates coordination overhead while preserving parallelism; no messaging needed.
    \item \textbf{Specialist agents with private, scoped memory} (Sections~\ref{sec:memory} and~\ref{sec:floor-amri}). Domain expertise accumulates across sessions, provided memory stays agent-private and canonical data lives in registries.
    \item \textbf{Adversarial skepticism as highest authority} (Section~\ref{sec:epistemic}). Counteracts the drift toward optimism when specialists evaluate their own domain.
    \item \textbf{Graded process rules} (Section~\ref{sec:floor-kcounter}). Rules earn mandatory status on accumulated evidence rather than by fiat, preventing both rule-rot and rule-inflation.
    \item \textbf{Context-free cross-axis verification} (Section~\ref{sec:floor-stage2}). A constructive, structurally sound answer to the multi-agent-agreement problem.
\end{enumerate}

\subsection{What Failed}

\begin{enumerate}
    \item \textbf{Large messaging teams} (more than $\sim$4 agents): notification avalanche, inbox contamination, agents completing tasks before receiving assignments. The lesson was learned repeatedly before the parallel-computation format was adopted.
    \item \textbf{Deferred recommendations}: items marked ``deferred'' were effectively lost; the carry-forward protocol was the fix, and required explicit enforcement.
    \item \textbf{Agent roleplaying}: agents writing sections designated for other agents produce confident, shallow output; prohibited by rule.
    \item \textbf{Comparing against derived rather than raw quantities}: the most expensive failure mode in the retraction log (Section~\ref{sec:assessment}).
    \item \textbf{The task-complete lie} (Section~\ref{sec:floor-closure}): agents claiming completion while a promised artifact is unwritten; addressed by on-disk completion verification and the mechanical-closure discipline.
    \item \textbf{Silent rule-rot} (Section~\ref{sec:floor-evoi}): an unconsumed priority table content-froze for roughly thirteen sessions while looking maintained; ``git-touched is not content-refreshed.''
    \item \textbf{Workshop inflation} (Section~\ref{sec:floor-workshop}): investigators padding workshop lists with carry-forward equations; addressed by the four-condition definition and three-question discriminator.
\end{enumerate}

\subsection{What Was Discovered About the Process Itself}

Two findings stand out. The first, present from the original draft: the free-form investigation mode---in which a coordinating agent directly explores the corpus making associative cross-domain connections---produced the project's signature early discovery (Section~\ref{sec:sacase}) and could not have been produced by any combination of specialists. The skill-based pipeline is valuable for \emph{execution} but can hinder \emph{discovery} where creative association is needed; the optimal methodology alternates between structured computation and unstructured exploration. The second, learned later: the alternation must be \emph{disciplined}, or the unstructured side degrades into bloviation (the workshop-inflation failure) and the structured side ossifies (the rule-rot failure). The methodology floor is, in large part, the machinery that keeps both sides honest.

% =====================================================================
\section{Discussion}
\label{sec:discussion}
% =====================================================================

\subsection{Generalizability}

The methodology was developed for a specific theoretical-physics project, but its core components---pre-registered gates with pinned machinery, source-authority hierarchies, scoped agent memory, working-paper coordination, carry-forward protocols, graded process rules, and context-free cross-axis verification---are domain-independent. Any research program involving multiple computation steps where later steps depend on earlier results, the possibility of negative results that constrain future work, several domains of expertise that must be combined, and a need for honest confidence calibration over time, could benefit from this epistemic architecture. The framing disciplines (Section~\ref{sec:floor-framing}) generalize as ``domain framing as auditable pre-registration'': fixing the direction of explanation and the sourcing of values turns a stylistic preference into a checkable predicate. Candidate domains include computational chemistry, genomics pipeline design, materials-science optimization, and any setting where an ensemble of agents must accumulate trustworthy results across sessions.

\subsection{The Role of the Human Researcher}

The human researcher served functions no agent took over. \textbf{Direction-setting}: deciding which questions to investigate, when to pivot, and when to stop; the project's paradigm shifts originated from the human, not from any agent. \textbf{Quality control}: rejecting agent output, correcting methodology, and enforcing epistemic discipline---the feedback-memory system records the explicit corrections that shaped subsequent sessions, several of which (the scope of memory under AMRI; the precise definition of a workshop) became standing rules. \textbf{Associative leaps and the publish decision}: certain cross-domain insights, and the decision of what to publish, remained human. The AI agents provided computational throughput, domain expertise across roughly three dozen specialist perspectives, systematic exploration, and queryable institutional memory. The combination---human direction with AI execution under a shared epistemic architecture---proved more productive than either alone.

\subsection{Limitations}

\begin{enumerate}
    \item \textbf{Context window}: each session starts with a finite context window that cannot hold the full project state. Compressed memory and---increasingly---queryable registries and the knowledge database (Section~\ref{sec:knowledge}) provide state, but the registry-promotion discipline (Section~\ref{sec:floor-amri}) only partly addresses the underlying compression loss.
    \item \textbf{Agent consistency}: the same agent type may behave differently across sessions due to stochastic sampling and to changes in the underlying model generation (the project spanned three; see below). Persistent memory mitigates but does not eliminate this.
    \item \textbf{Confirmation bias in closure}: agents sometimes compute what they already expect; the pre-registered-gate system limits but does not eliminate this tendency, and a fraction of computational effort goes to predictable closures (which nonetheless build the structural-wall catalogue).
    \item \textbf{Scalability}: the project's parallelism peaked at roughly two dozen agents in a single wave; whether the methodology scales to hundreds of agents is untested.
    \item \textbf{Model-generation drift}: the project ran across the Opus 4.6, 4.7, and 4.8 generations of the underlying model (the orchestrator's generation is recorded in session metadata---for example, an ``opus-4.7'' marker at session~81---and the current generation is 4.8). Attributing the work honestly requires acknowledging that the ``agents'' were not a fixed system but a succession of model generations under a stable epistemic architecture---which, if anything, strengthens the paper's central thesis that the architecture, not the agent, is the invariant.
\end{enumerate}

% =====================================================================
\section{Conclusion}
\label{sec:conclusion}
% =====================================================================

We have described an epistemic architecture for orchestrated AI research developed and refined over more than one hundred sessions of computational theoretical physics. The architecture's principal innovation is the separation of \emph{epistemic infrastructure} (rules, classifications, authority hierarchies, audit invariants, memory systems) from \emph{computational execution} (agent spawning, script running, data analysis). The infrastructure persists across sessions and accumulates institutional knowledge; the execution is ephemeral and parallelizable. The original draft of this paper described the first third of the project; the larger and, we now think, more durable contribution is the \emph{methodology floor}---the maturation of pre-registration into a formal under-specification taxonomy, of process rules into a graded promotion contract, and of multi-agent agreement into a context-free cross-axis verification protocol---that postdated that draft.

The project produced a curated set of standalone mathematical results that survive regardless of the framework's cosmological fate, a large and well-categorized retraction log that evidences a working falsification discipline, and a probability history that moved in both directions before the project deliberately suspended single-number probability reporting in favor of describing a constraint-surface geometry. That suspension is itself a methodological result: a discipline mature enough to recognize when a scalar would mislead.

The most important lesson is also the simplest: \emph{the rules matter more than the agents}. The same specialist agents---indeed, three successive generations of the underlying model---operating under different epistemic rules would produce different results: more optimistic, less rigorous, more prone to confirmation bias. The epistemic architecture is the invariant; the agents are the variables. The primary challenge of AI-assisted research is therefore not building better agents but designing better epistemic frameworks for them to operate within.

% =====================================================================
% References
% =====================================================================
\bibliographystyle{plain}

\begin{thebibliography}{99}

\bibitem{alphageometry}
    T. H. Trinh, Y. Wu, Q. V. Le, H. He, and T. Luong, ``Solving olympiad geometry without human demonstrations,'' \emph{Nature}, vol.~625, pp.~476--482, 2024.

\bibitem{alphafold}
    J. Jumper et al., ``Highly accurate protein structure prediction with AlphaFold,'' \emph{Nature}, vol.~596, pp.~583--589, 2021.

\bibitem{gnome}
    A. Merchant et al., ``Scaling deep learning for materials discovery,'' \emph{Nature}, vol.~624, pp.~80--85, 2023.

\bibitem{framework_paper}
    R. Berry, ``Phonon-exflation cosmology: spectral geometry of Jensen-deformed $SU(3)$ and the emergence of Standard Model and cosmological observables,'' companion manuscript, \emph{in preparation}, 2026.

\end{thebibliography}

\end{document}
